% ===========================================================================
% ΚΕΝΤΡΙΚΕΣ ΡΥΘΜΙΣΕΙΣ ΒΙΒΛΙΟΥ — ΞΕΚΙΝΑ ΑΠΟ ΕΔΩ
% ===========================================================================
% Αυτό είναι το πιο χρήσιμο αρχείο αν θέλεις να αλλάξεις την εμφάνιση μόνος σου.
% Τις τιμές τις αλλάζεις εδώ και όχι χειροκίνητα μέσα σε κάθε κεφάλαιο.
% Παράδειγμα: άλλαξε \BookLevel από Α΄ σε Β΄ Λυκείου.
% Κάθε γραμμή που αρχίζει με % είναι σχόλιο: δεν τυπώνεται στο βιβλίο.
% \newcommand{\Ονομα}{τιμή} ορίζει μια μεταβλητή για το υπόλοιπο έργο.
% ===========================================================================
% 1. Ταυτότητα βιβλίου / μεταδεδομένα

\newcommand{\BookSubject}{Φυσική}
\newcommand{\BookLevel}{Α΄ Λυκείου}
\newcommand{\BookTitle}{Φυσική Α΄ Λυκείου}
\newcommand{\BookSubtitle}{Θεωρία · Μεθοδολογία · Παραδείγματα · Ασκήσεις}
\newcommand{\BookPublisher}{ΕΚΔΟΤΙΚΟ ΠΡΟΤΥΠΟ}
\newcommand{\BookAuthors}{Δοκιμαστική μακέτα — όχι τελικό βιβλίο}
\newcommand{\BookEdition}{Καθολική μακέτα v0.32 · Οκτώβριος 2026}
\newcommand{\BookShortName}{Φυσική Α΄ Λυκείου}

% ---------------------------------------------------------------------------
% ΥΠΟΣΕΛΙΔΟ — ΔΥΟ ΑΝΕΞΑΡΤΗΤΕΣ ΡΥΘΜΙΣΕΙΣ, ΜΙΑ ΦΟΡΑ ΓΙΑ ΟΛΟ ΤΟ ΒΙΒΛΙΟ
% ---------------------------------------------------------------------------
% BookFooterMode:
%   full   = Φυσική Α΄ Λυκείου · σελ. 12   (διαθέσιμη εναλλακτική)
%   number = 12                           (ΠΡΟΕΠΙΛΟΓΗ: μόνο αριθμός)
% BookFooterPosition:
%   center = κεντραρισμένο σε κάθε σελίδα (εναλλακτική)
%   outer  = ΠΡΟΕΠΙΛΟΓΗ: αριστερά στις ζυγές, δεξιά στις μονές.
% Για αλλαγή τροποποίησε ΜΟΝΟ τις δύο παρακάτω τιμές.
% Ισχύουν αυτόματα για κανονικές, plain και εξεταστικές σελίδες.
% Εξώφυλλο, αρχή κεφαλαίου και σκόπιμα κενές σελίδες παραμένουν χωρίς υποσέλιδο.
% v0.24: χωρίς οριζόντια γραμμή κάτω, μόνο το περιεχόμενο (αριθμός).
% Η γραμμή της κεφαλίδας παραμένει αμετάβλητη.
\newcommand{\BookFooterMode}{number}
\newcommand{\BookFooterPosition}{outer}
% 2. Απαντήσεις στο PDF
% \ShowAnswerstrue: εμφανίζονται οι σύντομες απαντήσεις/υποδείξεις.
% \ShowAnswersfalse: αποκρύπτονται, χωρίς να σβήνεται το περιεχόμενό τους.

\newif\ifShowAnswers
\ShowAnswerstrue
% 3. Γραμματοσειρές — δεν διανέμουμε αρχεία γραμματοσειρών στο ZIP.
% Body: ελληνικό κυρίως κείμενο, Heading: τίτλοι, Math: εξισώσεις.
% Η STIX Two Math είναι ΠΡΟΣΩΡΙΝΗ δοκιμαστική επιλογή.
% Περιμένουμε την αισθητική έγκριση της Θάλειας.
% Για επιστροφή στην παλιά MathFontName: latinmodern-math.otf
% Για το βάρος Medium/Semibold αλλάζεις μόνο τα αντίστοιχα ονόματα εδώ.
% Μετά εκτελείς XeLaTeX δύο φορές για ενημέρωση παραπομπών και αρίθμησης.

\newcommand{\BodyFontName}{Inter}
\newcommand{\BodyBoldFontName}{Inter SemiBold}
\newcommand{\BodyItalicFontName}{Inter Italic}
\newcommand{\BodyBoldItalicFontName}{Inter SemiBold Italic}
\newcommand{\HeadingFontName}{Noto Serif}
\newcommand{\HeadingBoldFontName}{Noto Serif Medium}
\newcommand{\HeadingItalicFontName}{Noto Serif Italic}
\newcommand{\HeadingBoldItalicFontName}{Noto Serif Medium Italic}
\newcommand{\MathFontName}{STIXTwoMath-Regular.ttf}
% 4. ΕΙΚΟΝΙΔΙΑ ΔΙΔΑΚΤΙΚΩΝ ΚΟΥΤΙΩΝ — v0.25
% Η εμφάνιση Β εγκρίθηκε: μεγαλύτερα διανυσματικά εικονίδια στην αριστερή
% ζώνη, κεντραρισμένα στο ΥΨΟΣ του κουτιού, με κάθετη γραμμή αριστερά.
% Δεν υπάρχει πλέον σύμβολο μέσα στον τίτλο. Ο συγγραφέας εξακολουθεί
% να γράφει απλώς \begin{ThinkBox} ή \begin{AttentionBox}.
% \BookCalloutIconScale: κοινή κλίμακα για λάμπα και προειδοποιητικό τρίγωνο.
% \BookCalloutIconOffset: απόσταση του ΚΕΝΤΡΟΥ κάθε εικονιδίου από το
% αριστερό όριο του κουτιού. Το πλάτος της ζώνης ορίζεται στο κοινό
% στυλ BookCallout του styles/universalbook.sty.
\newcommand{\BookCalloutIconScale}{1.11}
% v0.34: η λάμπα είναι 10% μικρότερη από την κοινή κλίμακα.
% Η γεωμετρική κλίμακα δεν μειώνει τα πάχη γραμμής του TikZ.
\newcommand{\BookBulbScale}{0.90}
\newcommand{\BookBulbOutlineWidth}{1.08pt}
\newcommand{\BookBulbRayWidth}{.95pt}
\newcommand{\BookBulbFilamentWidth}{.85pt}
% v0.35: χαμηλό, πλατύ ανοιχτό βιβλίο με διακριτές καμπύλες σελίδων.
\newcommand{\BookOpenBookScale}{0.90}
\newcommand{\BookOpenBookOutlineWidth}{.90pt}
\newcommand{\BookOpenBookDetailWidth}{.68pt}
\newcommand{\BookCalloutIconOffset}{8.7mm}
% Ιστορική συμβατότητα: παλιές εντολές χαρακτήρων (όχι πλέον στους τίτλους).
% Διατηρούνται σε περίπτωση που χρησιμοποιούνται σε παλαιότερα αρχεία.
\newcommand{\SymbolFontName}{DejaVu Sans}
\newcommand{\AlertSymbolText}{▲}
\newcommand{\ThinkSymbolText}{◉}
\newcommand{\AlertSymbolSize}{9.0}
\newcommand{\ThinkSymbolSize}{9.0}
\newcommand{\AlertSymbolRaise}{0ex}
\newcommand{\ThinkSymbolRaise}{0ex}


% 5. Κείμενα ετικετών: ονόματα πλαισίων και ενοτήτων.
% Για γενικό βιβλίο π.χ. Μαθηματικών, τα αλλάζεις εδώ χωρίς επέμβαση στο sty.

% v0.29: ξεχωριστή ετικέτα μαθησιακών στόχων στην ΑΡΧΗ ΚΕΦΑΛΑΙΟΥ.
% Δεν επηρεάζει το περιβάλλον LearningGoals των φύλλων εργασίας.
\newcommand{\LabelChapterGoals}{Τι θα μάθουμε}
\newcommand{\LabelGoals}{Διδακτικοί στόχοι}
\newcommand{\LabelAlert}{ΠΡΟΣΟΧΗ}
\newcommand{\LabelThink}{Πώς σκέφτομαι}
\newcommand{\LabelHistory}{Ιστορίες από τη Φυσική}
\newcommand{\LabelMethod}{ΜΕΘΟΔΟΛΟΓΙΑ}
\newcommand{\LabelExample}{Παράδειγμα}
\newcommand{\LabelExercise}{Άσκηση}
\newcommand{\LabelSolution}{Λύση}
\newcommand{\LabelWorksheet}{ΦΥΛΛΟ ΕΡΓΑΣΙΑΣ}
% v0.15: ονομασίες των Κριτηρίων Αξιολόγησης — αλλάζουν κεντρικά.
\newcommand{\LabelAssessment}{ΚΡΙΤΗΡΙΟ ΑΞΙΟΛΟΓΗΣΗΣ}
\newcommand{\LabelAssessmentTopic}{ΘΕΜΑ}
\newcommand{\LabelAssessmentPoints}{Μονάδες}
\newcommand{\BookLogoFile}{} % e.g. assets/publisher-logo.pdf; leave empty for text logo
\newcommand{\LabelDefinition}{Ορισμός}
\newcommand{\LabelTheorem}{Θεώρημα}
\newcommand{\LabelProof}{Απόδειξη}

% Version v0.3 based on preferred typography B (Inter body, Noto Serif headings).
% Geometric Unicode label symbols: ▲ ΠΡΟΣΟΧΗ / ◉ Πώς σκέφτομαι.
% Font names are referenced by name; no font files are distributed.

% 6. Αριθμήσεις, κεφαλίδα και λεζάντες.
% Τα μεγέθη τυπογραφίας δίνονται σε pt εκτός αν γράφονται με μονάδα mm.
% \BookExerciseNumberSize: μέγεθος του μπλε «1 |» των άλυτων.
% \BookHeaderIconHeight: ύψος πουλιού στην κεφαλίδα και των δύο αντικριστών.

% 1 = show blue labels for all enumerate lists; actual numbers can be changed per list.
\newcommand{\BookListLabelColor}{BookSecondary}
\newcommand{\BookExerciseNumberSize}{11.4}
% Leave exercise titles empty in main problem sets: \begin{exercise}[ID]{}...
% Keep optional IDs in the source for future cross-references.
\newif\ifBookShowExerciseIDs
\BookShowExerciseIDsfalse
% Height of the publisher/header emblem; one value changes the entire book.
\newcommand{\BookHeaderIconHeight}{6.4mm}

% ---------------------------------------------------------------------------
% v0.22 — ΚΕΦΑΛΙΔΕΣ: κοινή μορφοποίηση και πεζοκεφαλαία ως προεπιλογή.
% ---------------------------------------------------------------------------
% Το κείμενο της ΚΕΦΑΛΙΔΑΣ είναι ανεξάρτητο από τη λωρίδα τίτλου.
% Για παράδειγμα, η γραμμή στην κορυφή μπορεί να γράφει
% «Κριτήριο αξιολόγησης 1», ενώ η μπλε λωρίδα παραμένει
% «ΚΡΙΤΗΡΙΟ ΑΞΙΟΛΟΓΗΣΗΣ 1» από τη μεταβλητή \LabelAssessment.
% Η \nouppercase στο sty αποτρέπει την αυτόματη κεφαλαιοποίηση των τίτλων,
% ενώ εδώ δίνουμε απευθείας την επιθυμητή γραφή της ετικέτας.
\newcommand{\BookChapterHeaderPrefix}{Ενότητα}
\newcommand{\BookAssessmentHeaderLabel}{Κριτήριο αξιολόγησης}
% Κοινή γραμματοσειρά, μέγεθος και βάρος για τους δύο τύπους κεφαλίδας:
\newcommand{\BookHeaderLabelFont}{\sffamily\scriptsize\bfseries}
\newcommand{\BookHeaderTitleFont}{\sffamily\scriptsize}
% Χρωματικοί ρόλοι: αλλάζουν αυτομάτως μεταξύ SCREEN και PRINT.
\newcommand{\BookHeaderLabelColor}{BookPrimary}
\newcommand{\BookHeaderTitleColor}{BookInk}
\newcommand{\BookHeaderRuleColor}{BookSecondary}
% Το πλάτος γραμμής και η απόσταση τίτλου/πουλιού ορίζονται μια φορά.
\newcommand{\BookHeaderRuleWidth}{.35pt}
\newcommand{\BookHeaderIconGap}{1.2mm}
% Το ύψος του πουλιού ρυθμίζεται από την BookHeaderIconHeight παραπάνω.
% Η θέση του μένει εξωτερικά στις αντικριστές σελίδες, και στις ζυγές
% το σήμα κατοπτρίζεται χωρίς αλλαγή εικόνας.

% 6α. Αντικριστά περιθώρια: εσωτερικό = πλευρά ράχης, εξωτερικό = άκρη.
% Δοκιμαστικές τιμές, ΟΧΙ οριστικές προδιαγραφές τυπογραφείου.
% Εδώ 17 + 20 mm αντί για 20 + 16 mm, ώστε να μείνει περισσότερος χώρος έξω.
% Το class book με twoside αντιστρέφει αυτόματα τα περιθώρια σε ζυγές σελίδες.
\newcommand{\BookMarginInner}{17mm}
% Καθαρό μέγεθος δοκιμίου. Το ξάκρισμα παραμένει ανοικτή προδιαγραφή.
\newcommand{\BookTrimWidth}{170mm}
\newcommand{\BookTrimHeight}{240mm}
\newcommand{\BookMarginOuter}{20mm}
\newcommand{\BookMarginTop}{19mm}
\newcommand{\BookMarginBottom}{20mm}

% 6β. Η αρχή κάθε κεφαλαίου τοποθετείται πάντα σε ΜΟΝΗ/ΔΕΞΙΑ σελίδα.
% Αυτό υλοποιείται με την απλή \cleardoublepage στο universalbook.sty.
% Αν η επόμενη σελίδα θα ήταν αριστερή, το LaTeX προσθέτει ΜΙΑ λευκή.
% ΔΕΝ απαιτείται λευκή η αριστερή σελίδα απέναντι από κάθε τίτλο.
% Δεν υπάρχει πλέον πρόσθετη επιλογή ή ειδικός μετρητής για αυτό.

\newcommand{\BookExerciseSeparatorWidth}{0.75pt}
\newcommand{\BookExerciseSeparatorHeight}{11.0pt}
% Float caption settings are controlled here, not inside chapters.
\newcommand{\BookCaptionSize}{9.1}
\newcommand{\BookCaptionLineHeight}{12}
\newcommand{\BookFigureCaptionName}{Σχήμα}
\newcommand{\BookTableCaptionName}{Πίνακας}
% v0.23: κοινές παράμετροι ΚΑΙ για τους δύο πίνακες.
% Αλλάζεις τον χρωματισμό μόνο εδώ, όχι σε κάθε tabularx.
% PRINT: τα BookPrimary / BookRowTint μετατρέπονται αυτόματα σε
% 100% / 10% του προσωρινού δεύτερου μελανιού. Κανένα νέο μελάνι.
\newcommand{\BookTableHeaderFill}{BookPrimary}
\newcommand{\BookTableStripeFill}{BookRowTint}
\newcommand{\BookTableRowStretch}{1.33}
% When changed to false, figure/table descriptions become upright instead of italic.
\newif\ifBookItalicCaptions
\BookItalicCaptionstrue
\newcommand{\BookChapterPhotoHeight}{60mm}
% v0.29: λεπτή, διακριτική πλευρική γραμμή στο «Τι θα μάθουμε».
% BookSecondary = μεσαίο μπλε στην SCREEN / 55% του 2ου μελανιού PRINT.
% Καμία αλλαγή στο κείμενο των \item από τον συγγραφέα.
\newcommand{\BookOpeningGoalsAccent}{BookSecondary}
\newcommand{\BookOpeningGoalsRuleWidth}{1.35pt}
\newcommand{\BookOpeningGoalsLeftPadding}{4.3mm}

% 7. Ιεραρχία τίτλων και πλαίσια λυμένων παραδειγμάτων.
% Κάθε τιμή *true* ανοίγει μια οπτική επιλογή· *false* την απενεργοποιεί.
% Με BookExampleDarkHeadertrue ο τίτλος είναι λευκός σε σκούρα μπλε λωρίδα.
% Με BookMajorHeadingBandtrue οι βασικές ενότητες έχουν δική τους μπλε μπάρα.

% Example title in dark-blue band (true), or on pale-blue background with a hairline (false).
\newif\ifBookExampleDarkHeader
\BookExampleDarkHeadertrue
% Major sections: compact blue bar with white text (true), legacy heading + rule (false).
\newif\ifBookMajorHeadingBand
\BookMajorHeadingBandtrue
% Small subsection rules: solid thin rule (false), dotted leader (true).
\newif\ifBookQuestionDottedRules
\BookQuestionDottedRulesfalse
% Compact heading bands, distinct from the running header and worksheet banners.
\newcommand{\BookMajorBandSize}{11.7}
\newcommand{\BookMajorBandTopPad}{1.35mm}
\newcommand{\BookMajorBandBottomPad}{1.35mm}
% True/false checkboxes and spacing. Length is in centimeters.
\newcommand{\BookCheckSquareSize}{.27}
\newcommand{\BookTrueFalseGap}{3.1mm}

% 8. Μικροτυπογραφία: διάκενα, ύψη γραμμών και βάρη ετικετών.
% Χρήσιμο όταν ζητηθεί αλλαγή από την επιμελήτρια των εκδόσεων.

\newcommand{\BookExampleTitleFontSize}{10.4}
\newcommand{\BookExampleTitleLineHeight}{13.5}
\newcommand{\BookExampleBarPad}{2.0mm}
% Κοινό μέγεθος για τίτλους ThinkBox / AttentionBox / HistoryBox.
\newcommand{\BookCalloutLabelSize}{9.2}
\newcommand{\BookCalloutLabelLeading}{11.5}
% Παλιό όνομα για αναφορά· οι αλλαγές γίνονται στο BookCalloutLabelSize.
\newcommand{\BookThinkLabelSize}{\BookCalloutLabelSize}
\newcommand{\BookSolutionLabelSize}{10.2}
\newcommand{\BookTheoryRuleWidth}{0.55pt}
\newcommand{\BookQuestionRuleWidth}{0.35pt}

\newcommand{\BookMediumFontName}{Inter Medium}

% ==========================================================================
% 9. ΓΡΑΦΙΚΑ — ΚΕΝΤΡΙΚΑ ΜΕΓΕΘΗ (v0.14)
% Τα παρακάτω χρησιμοποιούνται μόνο στα κοινά στυλ TikZ/PGFPlots.
% Δεν χρειάζεται να γράφεις αριθμούς pt μέσα σε κάθε διάγραμμα.
% Μεγάλωσε βέλη/σύμβολα εδώ και θα ενημερωθούν όλα τα νέα σχήματα.
% ==========================================================================
\newcommand{\BookFigAxisWidth}{0.72pt}    % κύριοι άξονες
\newcommand{\BookFigVectorWidth}{1.12pt}  % βέλη φυσικών μεγεθών
\newcommand{\BookFigGraphWidth}{1.05pt}   % κύριες καμπύλες
\newcommand{\BookFigAuxWidth}{0.45pt}     % βοηθητικά/οδηγοί
\newcommand{\BookFigArrowLength}{2.5mm}  % μήκος αιχμής
\newcommand{\BookFigArrowWidth}{1.7mm}   % άνοιγμα αιχμής
